\chapter{Connecting a PLC to the Adapter}\label{ch:plc}

This chapter tells you how to make an EDS file for your adapter and how to
connect a PLC to it. Do the checks in section~\ref{sec:precheck} before you
start the PLC software.

\section{Checks Before You Use the PLC}\label{sec:precheck}

\begin{procedure}{Check the adapter from a PC}
\begin{enumerate}
  \item Connect a PC to the same network as the adapter.
  \item Build \code{eiptool} on the PC (section~\ref{sec:hostbuild}).
  \item Read the identity of the adapter:
\begin{lstlisting}[language=shell]
eiptool identity <adapter-ip>
\end{lstlisting}
  \item Make sure that the vendor, product code, and name are correct.
  \item Read the data of the produced assembly:
\begin{lstlisting}[language=shell]
eiptool get <adapter-ip> 4 100 3
\end{lstlisting}
  \item Make sure that the byte count is equal to the assembly size.
\end{enumerate}
\end{procedure}

\section{Make the EDS File}\label{sec:eds}

The file \file{examples/eds/QuickEIP_Adapter.eds} agrees with the assemblies
in section~\ref{sec:start}. Use it as the start for your EDS file.
Table~\ref{tab:edsfields} shows which values must agree with your
configuration.

\begin{table}[htbp]
\centering
\caption{EDS fields and the configuration values that they must agree with}\label{tab:edsfields}
\small
\begin{tabularx}{\linewidth}{@{}P{0.38\linewidth}L@{}}
\toprule
\textbf{EDS field} & \textbf{Must be equal to} \\
\midrule
\code{VendCode}, \code{ProdType}, \code{ProdCode} & \code{vendor_id}, \code{device_type}, \code{product_code} \\
\code{MajRev}, \code{MinRev} & \code{revision_major}, \code{revision_minor}. EDS tools need 1--255, so do not use a minor revision of 0. \\
\code{ProdName} & \code{product_name} \\
\code{VendName}, \code{Catalog} & Your company and your part number. No configuration field agrees with these. \\
\code{Assem}\textit{NNN}, its size, and its member size in bits & An \code{EipAssemblyConfig} instance, its \code{size}, and \code{size} $\times$ 8 \\
Connection sizes, \code{TSpec1}, \code{TSpec2} & The assembly sizes (not the wire sizes) \\
\code{Number_Of_Static_Instances} under \code{[Assembly]} & \code{assembly_count} \\
Connection path & \code{20 04 24} \textit{config} \code{2C} \textit{consumed} \code{2C} \textit{produced}, with instances in hexadecimal, for instances up to 255 \\
Assembly path & \code{20 04 24} \textit{instance} \code{30 03}, with the instance in hexadecimal \\
\bottomrule
\end{tabularx}
\end{table}

For the example assemblies, config 151 is \hex{97}, consumed 150 is \hex{96},
and produced 100 is \hex{64}. Thus the connection path is
\code{20 04 24 97 2C 96 2C 64}.

\begin{procedure}{Adapt the sample EDS file}
\begin{enumerate}
  \item Copy \file{examples/eds/QuickEIP_Adapter.eds} and give the copy a new name.
  \item Set the \code{[Device]} values from table~\ref{tab:edsfields}.
  \item Set each \code{Assem}\textit{NNN} entry to agree with one assembly.
  \item Set the connection sizes and \code{TSpec1} and \code{TSpec2} to the assembly sizes.
  \item Set the connection path to your config, consumed, and produced instances.
  \item Do not change the other values.
  \item Do the checks of section~\ref{sec:precheck} again.
\end{enumerate}
\end{procedure}

\begin{caution}
When an assembly layout changes, also change \code{ProdCode} or the revision.
If you do not, a PLC can use an old EDS file with the new device, and the
data has a different meaning.
\end{caution}

\begin{note}
\notelabel The connection values in the sample file were checked: they opened
a connection to a running adapter. But nobody has imported the file itself
into PLC software yet.
\end{note}

Appendix~\ref{app:tools} contains the important parts of the sample file.

\section{Micro850 (Connected Components Workbench)}

This setup was tested with the ClearCore port.

\begin{procedure}{Connect a Micro850}
\begin{enumerate}
  \item In Connected Components Workbench, add a generic EtherNet/IP device.
  \item Set its address to the address of the adapter.
  \item Set the output assembly (\OT) to instance 150, 32 bytes.
  \item Set the input assembly (\TO) to instance 100, 32 bytes.
  \item Set the config assembly to instance 151, 0 bytes.
  \item Download the project to the controller.
  \item Put the controller in Run mode.
  \item On the adapter, make sure that \code{io_open} and \code{run} are 1 in
        the snapshot.
\end{enumerate}
\end{procedure}

Table~\ref{tab:micro850} shows the values that a Micro850 was seen to send.

\begin{table}[htbp]
\centering
\caption{Connection values sent by a Micro850}\label{tab:micro850}
\begin{tabularx}{\linewidth}{@{}P{0.32\linewidth}L@{}}
\toprule
\textbf{Setting} & \textbf{Value} \\
\midrule
RPI & 50 ms in each direction \\
\OT & Point-to-point, with the run/idle header \\
\TO & No header \\
Timeout multiplier & 0 ($\times$4, thus 200 ms at a 50 ms RPI) \\
Electronic key & All zeros \\
Config & Names the config instance and sends no data \\
\bottomrule
\end{tabularx}
\end{table}

The config instance must exist. If it does not, the adapter refuses the open
with \hex{01} / \hex{0129}. A size-0 \code{EIP_ASM_CONFIG} assembly is
sufficient.

\section{CompactLogix and ControlLogix (Studio 5000)}

This adapter has not been tested with these controllers. But an EDS file with
the same connection values, and generic module settings in this form, have
operated on a CompactLogix with a different adapter.

You can install the EDS file with the \menu{EDS Hardware Installation Tool}
and add the device as a module. Or you can add a \menu{Generic Ethernet
Module} with the values in table~\ref{tab:generic}.

\begin{table}[htbp]
\centering
\caption{Generic Ethernet Module settings}\label{tab:generic}
\begin{tabularx}{\linewidth}{@{}P{0.32\linewidth}L@{}}
\toprule
\textbf{Field} & \textbf{Value} \\
\midrule
\menu{Comm Format} & Data -- SINT \\
\menu{IP Address} & The address of the adapter \\
\menu{Input} & Assembly instance 100, size 32 \\
\menu{Output} & Assembly instance 150, size 32 \\
\menu{Configuration} & Assembly instance 151, size 0 \\
\bottomrule
\end{tabularx}
\end{table}

Then set the RPI on the \menu{Connection} tab. In the two methods, the config
assembly must exist on the adapter, also at size 0. It comes first in the
connection path.

\section{Adapter on a PC (POSIX or Windows Port)}

This setup has not been tested with a PLC.

\begin{procedure}{Prepare a PC adapter for a PLC}
\begin{enumerate}
  \item In the firewall, open TCP 44818, UDP 44818, and UDP 2222.
  \item If the PC has more than one interface, set \code{EIP_NETIF} to the
        interface that the PLC uses.
  \item Under WSL2, use mirrored networking.
  \item Set the RPI in the PLC to 50 ms or more.
  \item Set \code{timeout_floor_ms} to 500 or more.
\end{enumerate}
\end{procedure}

A PC has less regular timing than an embedded target. Thus start with a
longer RPI.

\section{Rules for All PLCs}

\begin{itemize}
  \item The adapter accepts one Class~1 connection at a time. It refuses a
        second originator.
  \item Keep the RPI much longer than the clock step of your port. On a
        NetBurner at 20 Hz, use 50 ms or more (section~\ref{sec:netburner}).
  \item If the PLC shows a fault, find the status in
        table~\ref{tab:forwardopen}. \code{16\#0127} and \code{16\#0128} are
        size errors.
  \item Outputs read as zeros while the PLC is in Program mode.
\end{itemize}
